\documentclass[11pt, a4paper]{article}

% --- Paquetes básicos ---
\usepackage[utf8]{inputenc}
\usepackage[T1]{fontenc}
\usepackage[spanish]{babel}
\usepackage{graphicx} % Para insertar imágenes
\usepackage{lipsum}   % Para generar texto de relleno
\usepackage{amsmath}  % Para fórmulas matemáticas
\usepackage{hyperref} % Para enlaces clicables en el índice y referencias

% --- Metadatos del documento ---
\title{Mi Primer Documento en \LaTeX}
\author{Tu Nombre}
\date{\today}

\begin{document}

\maketitle

\begin{abstract}
  Este es un documento mínimo de prueba («Hello World») que demuestra la estructura básica de un artículo en \LaTeX: secciones, texto simulado, ecuaciones, una figura estándar y referencias bibliográficas.
\end{abstract}

\section{Introducción}
\label{sec:intro}

Bienvenido a \LaTeX. A diferencia de un procesador de textos tradicional, aquí escribes la estructura lógica del documento y el motor tipográfico se encarga del diseño. Por ejemplo, podemos escribir ecuaciones en línea como $E = mc^2$, o centradas y numeradas:

\begin{equation}
  \int_{-\infty}^{\infty} e^{-x^2} \, dx = \sqrt{\pi}
  \label{eq:gauss}
\end{equation}

Como se observa en la Ecuación~\eqref{eq:gauss} de la Sección~\ref{sec:intro}, \LaTeX{} gestiona automáticamente la numeración y las referencias cruzadas.

\section{Texto de ejemplo y figuras}

A continuación se muestra un párrafo generado automáticamente mediante el paquete \texttt{lipsum}, junto con una referencia al manual clásico de \LaTeX~\cite{lamport1994}:

\lipsum[1]

En la Figura~\ref{fig:ejemplo} se incluye una imagen de prueba incluida por defecto en Overleaf, por lo que no necesitas subir ningún archivo adicional para compilar este ejemplo:

\begin{figure}[htbp]
  \centering
  \includegraphics[width=0.55\textwidth]{example-image-a}
  \caption{Imagen de ejemplo integrada en la distribución TeX Live.}
  \label{fig:ejemplo}
\end{figure}

% --- Bibliografía autocontenida ---
\begin{thebibliography}{9}

  \bibitem{lamport1994}
  Leslie Lamport,
  \textit{\LaTeX: A Document Preparation System},
  2.\textsuperscript{a} edición, Addison-Wesley, 1994.

\end{thebibliography}

\end{document}